\section{THỰC TRẠNG HOẠT ĐỘNG QUẢN TRỊ TẠI CÔNG TY CỔ PHẦN FPT}

\subsection{Khái niệm và sơ đồ cơ cấu quản trị công ty}
Theo khoa học quản lý hiện đại, \textbf{Quản trị công ty (Corporate Governance)} là hệ thống các cơ chế, quy trình, mối quan hệ và nguyên tắc mà thông qua đó doanh nghiệp được định hướng chiến lược và kiểm soát. Quản trị công ty xác định cơ cấu phân bổ quyền hạn và trách nhiệm giữa các chủ thể chính: Đại hội đồng Cổ đông, Hội đồng Quản trị, Ban Kiểm soát và Ban Điều hành.

Sơ đồ tổ chức cơ cấu quản trị tại FPT được thiết lập chặt chẽ theo Luật Doanh nghiệp 2020 và các thông lệ quản trị đại chúng tiên tiến:

\begin{center}
\begin{verbatim}
               ĐẠI HỘI ĐỒNG CỔ ĐÔNG (ĐHĐCĐ)
                            |
            +---------------+---------------+
            |                               |
    BAN KIỂM SOÁT                 HỘI ĐỒNG QUẢN TRỊ (HĐQT)
    (Độc lập giám sát)           (7 Thành viên, 3 độc lập)
                                            |
                         +------------------+------------------+
                         | Các Ủy ban & Ban trực thuộc HĐQT:      |
                         | - Ủy ban Nhân sự và Lương thưởng       |
                         | - Ủy ban Chính sách Phát triển         |
                         | - Văn phòng Chủ tịch HĐQT              |
                         | - Bộ phận Kiểm toán Nội bộ             |
                         +------------------+------------------+
                                            |
                                  BAN ĐIỀU HÀNH (CEO)
\end{verbatim}
\end{center}

\subsection{Đại hội đồng Cổ đông (ĐHĐCĐ)}
Đại hội đồng Cổ đông là cơ quan có thẩm quyền quyết định cao nhất của Công ty Cổ phần FPT, bao gồm tất cả cổ đông có quyền biểu quyết (tổng số cổ phần lưu hành sau đợt chia thưởng 10:1 năm 2026 đạt 1.885.727.648 cổ phần). ĐHĐCĐ thực hiện quyền năng thông qua định hướng phát triển chiến lược, phê duyệt báo cáo tài chính kiểm toán hàng năm, phân phối lợi nhuận, mức chi trả cổ tức, bầu/miễn nhiệm thành viên HĐQT và Ban Kiểm soát.

Đặc biệt, trong các kỳ họp ĐHĐCĐ (bao gồm năm 2025), FPT đã tiên phong ứng dụng công nghệ Trí tuệ Nhân tạo (AI) trên nền tảng FPT.AI trong việc nhận diện sinh trắc học khuôn mặt và tự động đối soát thông tin định danh của cổ đông, rút ngắn thời gian xác thực của mỗi cổ đông xuống chỉ còn \textbf{khoảng 2 giây}, bảo đảm tính chính xác, minh bạch và an toàn.

\subsection{Hội đồng Quản trị (HĐQT) và các nhân vật lãnh đạo}
Hội đồng Quản trị FPT là cơ quan đại diện cho cổ đông, có toàn quyền nhân danh công ty để quyết định mọi vấn đề liên quan đến mục tiêu và lợi ích của công ty, ngoại trừ những vấn đề thuộc thẩm quyền của ĐHĐCĐ.

\subsubsection{Cơ cấu nhân sự HĐQT nhiệm kỳ 2022--2027}
HĐQT FPT nhiệm kỳ 2022--2027 gồm \textbf{7 thành viên}, trong đó có 3 thành viên độc lập (chiếm tỷ lệ 42,8\% -- vượt mức tối thiểu theo luật định), bảo đảm tính đa dạng về chuyên môn, quốc tịch và tính khách quan trong việc ra quyết định (Bảng \ref{tab:hdqt_fpt}):

\begin{table}[htbp]
\centering
\caption{Danh sách thành viên Hội đồng Quản trị FPT nhiệm kỳ 2022--2027}
\label{tab:hdqt_fpt}
\begin{tabularx}{\textwidth}{p{4.2cm} p{4.5cm} X}
\toprule
\textbf{Họ và tên} & \textbf{Chức danh tại HĐQT} & \textbf{Vai trò chuyên môn} \\
\midrule
\textbf{Ông Trương Gia Bình} & Chủ tịch HĐQT & Sáng lập viên FPT, định hình tầm nhìn chiến lược, ngoại giao cấp cao \\
\textbf{Ông Bùi Quang Ngọc} & Phó Chủ tịch HĐQT & Đồng sáng lập, chuyên gia quản trị hệ thống và kiểm soát vận hành \\
\textbf{Ông Đỗ Cao Bảo} & Ủy viên HĐQT & Đồng sáng lập, phụ trách chính sách phát triển nhân tài \\
\textbf{Ông Jean-Charles Belliol} & Ủy viên HĐQT độc lập & Chuyên gia tài chính, đầu tư quốc tế (quốc tịch Pháp) \\
\textbf{Ông Hiroshi Yokotsuka} & Ủy viên HĐQT độc lập (đến 4/2026) & Nguyên Chủ tịch Hiệp hội CNTT Nhật Bản (quốc tịch Nhật Bản); miễn nhiệm 16/4/2026 theo nguyện vọng cá nhân \\
\textbf{Ông Toshikazu Nambu} & Ủy viên HĐQT (từ 4/2026) & Phó Chủ tịch Tập đoàn Sumitomo (Nhật Bản), được bầu bổ sung ngày 16/4/2026 thay ông Yokotsuka \\
\textbf{Ông Hampapur R. Binod} & Ủy viên HĐQT độc lập & Cựu lãnh đạo cấp cao Tập đoàn công nghệ Infosys (quốc tịch Ấn Độ) \\
\textbf{Bà Trần Thị Hồng Lĩnh} & Ủy viên HĐQT & Đại diện phần vốn nhà nước (SCIC) \\
\bottomrule
\end{tabularx}
\end{table}

Trong năm 2025, HĐQT FPT đã tổ chức tổng cộng \textbf{7 cuộc họp} (gồm 4 cuộc họp trực tiếp và 3 lần lấy ý kiến bằng văn bản) với tỷ lệ tham dự đạt xấp xỉ 100\%, quyết định các chủ trương lớn về ngân sách, đầu tư M\&A và chiến lược phát triển công nghệ AI. Tại ĐHĐCĐ thường niên ngày 16/4/2026, cổ đông đã thông qua miễn nhiệm ông Hiroshi Yokotsuka và bầu bổ sung ông Toshikazu Nambu vào HĐQT nhiệm kỳ 2022--2027.

\subsubsection{Các cơ quan và ủy ban chuyên trách trực thuộc HĐQT}
Nhằm nâng cao chất lượng quản trị và tính chuyên môn hóa, HĐQT FPT duy trì 4 cơ quan trực thuộc:
\begin{enumerate}
    \item \textbf{Ủy ban Nhân sự và Lương thưởng:} Do ông Đỗ Cao Bảo chủ trì; tham mưu chính sách nhân sự cấp cao, cơ chế đãi ngộ, quy hoạch bổ nhiệm và đánh giá hiệu quả làm việc của Ban Điều hành.
    \item \textbf{Ủy ban Chính sách Phát triển:} Do ông Trương Gia Bình làm Chủ tịch; nghiên cứu xu thế công nghệ thế giới, hoạch định và giám sát việc thực thi chiến lược chuyển đổi số, bán dẫn và AI.
    \item \textbf{Văn phòng Chủ tịch HĐQT:} Do bà Phạm Ngọc Anh làm Chánh Văn phòng; điều phối công tác hành chính quản trị, đối ngoại cổ đông và chuẩn bị tài liệu nghị quyết.
    \item \textbf{Bộ phận Kiểm toán Nội bộ:} Thành lập từ năm 2022, trực thuộc trực tiếp HĐQT; đánh giá độc lập tính hữu hiệu của hệ thống kiểm soát nội bộ, quản trị rủi ro và giám sát tuân thủ.
\end{enumerate}

\subsection{Ban Kiểm soát (BKS)}
Ban Kiểm soát hoạt động độc lập với HĐQT và Ban Điều hành công ty mẹ, do ĐHĐCĐ trực tiếp bầu ra nhằm bảo vệ quyền và lợi ích hợp pháp của các cổ đông:
\begin{itemize}
    \item \textbf{Ông Nguyễn Việt Thắng:} Trưởng Ban Kiểm soát (gia nhập FPT từ năm 1992, Phó Hiệu trưởng ĐH FPT kiêm Viện trưởng Viện FSB).
    \item \textbf{Bà Dương Thùy Dương:} Thành viên Ban Kiểm soát (từ năm 2022).
    \item \textbf{Ông Nguyễn Khải Hoàn:} Thành viên Ban Kiểm soát (đồng thời đảm nhiệm vai trò Phó TGĐ kiêm COO tại công ty con FPT Software và CEO FPT Americas từ 01/2026; không kiêm nhiệm trong HĐQT hay Ban Điều hành của công ty mẹ FPT).
\end{itemize}
Ban Kiểm soát thực hiện giám sát tính hợp pháp, hợp lý và trung thực trong việc quản lý, điều hành hoạt động kinh doanh; kiểm tra độc lập các báo cáo tài chính quý, bán niên và năm trước khi công bố ra thị trường.

\subsection{Người đại diện theo pháp luật}
Tại cấp Công ty Cổ phần FPT (Công ty mẹ), người đại diện theo pháp luật được ghi nhận chính thức là \textbf{Ông Trương Gia Bình} (Chủ tịch HĐQT). Ngoài ra, Điều lệ công ty quy định Tổng Giám đốc (\textbf{Ông Nguyễn Văn Khoa}) đồng thời là người đại diện theo pháp luật trong các phạm vi điều hành tác nghiệp được phân quyền.

Tại các đơn vị cấp dưới, người đại diện theo pháp luật được bổ nhiệm nhằm phân định rõ trách nhiệm pháp lý:
\begin{itemize}[noitemsep]
    \item \textit{FPT IS:} Ông Nguyễn Hoàng Minh và Ông Trần Đăng Hòa.
    \item \textit{FPT Smart Cloud:} Ông Lê Hồng Việt (Tổng Giám đốc).
    \item \textit{FPT Telecom:} Ông Hoàng Việt Anh (Chủ tịch) và Ông Nguyễn Hoàng Linh (Tổng Giám đốc).
    \item \textit{FPT Retail:} Ông Hoàng Trung Kiên (Tổng Giám đốc).
\end{itemize}

\subsection{Phân định ranh giới giữa Quản trị (Governance) và Quản lý (Management)}
Một điểm sáng quan trọng trong lịch sử quản trị của FPT là việc \textbf{tách bạch vị trí Chủ tịch HĐQT và Tổng Giám đốc từ năm 2013}. Động thái này giúp loại bỏ sự tập trung quyền lực tuyệt đối, tránh hiện tượng ``vừa đá bóng, vừa thổi còi'', đảm bảo HĐQT thực hiện chức năng giám sát khách quan đối với Ban Điều hành.

Sự phân định ranh giới giữa Quản trị và Quản lý tại FPT được thể hiện rõ qua Bảng \ref{tab:so_sanh_quan_tri_quan_ly}:

\begin{table}[htbp]
\centering
\caption{Phân định ranh giới giữa Quản trị và Quản lý tại Công ty Cổ phần FPT}
\label{tab:so_sanh_quan_tri_quan_ly}
\begin{tabularx}{\textwidth}{p{3.2cm} X X}
\toprule
\textbf{Tiêu chí} & \textbf{Quản trị (Governance)} & \textbf{Quản lý (Management)} \\
\midrule
\textbf{Chủ thể thực hiện} & Đại hội đồng Cổ đông, Hội đồng Quản trị & Ban Điều hành (Tổng Giám đốc, các Phó TGĐ, Giám đốc khối) \\
\textbf{Bản chất hoạt động} & Thiết lập tầm nhìn, phê duyệt chiến lược, ban hành chính sách và kiểm soát & Thực thi chiến lược, tổ chức sản xuất kinh doanh, vận hành thường nhật \\
\textbf{Trọng tâm câu hỏi} & \textit{``Làm cái gì?''} (What) và \textit{``Vì ai?''} (For Whom) & \textit{``Làm như thế nào?''} (How) và \textit{``Ai làm?''} (Who) \\
\textbf{Tầm nhìn thời gian} & Dài hạn (3--5 năm, 10 năm) & Ngắn và trung hạn (hàng ngày, tháng, quý, năm tài chính) \\
\textbf{Mối quan hệ quyền lực} & Đại diện chủ sở hữu, giám sát bộ máy điều hành & Chịu sự bổ nhiệm, đánh giá và giám sát trực tiếp từ HĐQT \\
\textbf{Trách nhiệm giải trình} & Chịu trách nhiệm trước toàn thể cổ đông và pháp luật & Chịu trách nhiệm trực tiếp trước Hội đồng Quản trị \\
\bottomrule
\end{tabularx}
\end{table}

\subsection{Đối chiếu với các chuẩn mực quản trị quốc tế}
FPT là một trong những doanh nghiệp niêm yết áp dụng sớm nhất \textit{Bộ nguyên tắc Quản trị Công ty theo thông lệ tốt nhất tại Việt Nam (Vietnam CG Code)} và các khuyến nghị của OECD:
\begin{itemize}
    \item \textbf{Tính độc lập và đa dạng:} Tỷ lệ thành viên HĐQT độc lập đạt 3/7, hội tụ các chuyên gia quốc tế hàng đầu từ Nhật Bản, Pháp, Ấn Độ.
    \item \textbf{Tính minh bạch thông tin:} Công bố thông tin tài chính song ngữ Việt -- Anh kịp thời, tổ chức các buổi gặp gỡ nhà đầu tư định kỳ hàng quý (Investor Relations).
    \item \textbf{Cơ chế kiểm soát nội bộ:} Thành lập bộ phận Kiểm toán Nội bộ độc lập trực thuộc HĐQT, ngăn ngừa xung đột lợi ích giữa công ty mẹ và các công ty liên kết.
\end{itemize}
